\medskip

\textbf{PARTIE A}

\medskip

On considère l'équation différentielle 
\[(E) : \quad y' +\dfrac14y = 20\e^{-\frac14 x},\]
d'inconnue $y$, fonction définie et dérivable sur l'intervalle $[0~;~ +\infty[$.

\medskip

\begin{enumerate}
\item Déterminer la valeur du réel $a$ tel que la fonction $g$ définie sur l'intervalle $[0~;~ +\infty[$ par $g(x) = ax\e^{-\frac14 x}$ soit une solution particulière de l'équation différentielle $(E)$.

\item On considère l'équation différentielle 
\[(E') : \quad y' +\dfrac14 y = 0,\]
d'inconnue $y$, fonction définie et dérivable sur l'intervalle $[0~;~ +\infty[$.

Déterminer les solutions de l'équation différentielle $(E')$.
\item En déduire les solutions de l'équation différentielle $(E)$.
\item Déterminer la solution $f$ de l'équation différentielle $(E)$ telle que $f(0) = 8$.
\end{enumerate}

\medskip

\textbf{PARTIE B}

\medskip

On considère la fonction $f$ définie sur l'intervalle $[0~;~ +\infty[$ par 
\[f(x) = (20x + 8)\e^{-\frac14 x}.\]

On admet que la fonction $f$ est dérivable sur l'intervalle $[0~;~ +\infty[$ et on note $f'$ sa fonction dérivée sur l'intervalle $[0~;~ +\infty[$. De plus, on admet que $\displaystyle\lim_{x \to + \infty}f(x) = 0$.

\medskip

\begin{enumerate}
\item 
	\begin{enumerate}
		\item Justifier que, pour tout réel $x$ positif, 
\[f'(x) = (18 - 5x)\e^{-\frac14 x}.\]

		\item En déduire le tableau de variations de la fonction $f$. On précisera la valeur exacte du maximum de la fonction $f$ sur l'intervalle $[0~;~ +\infty[$.
	\end{enumerate}
\item Dans cette question on s'intéresse à l'équation $f(x) = 8$.
	\begin{enumerate}
		\item Justifier que l'équation $f(x) = 8$ admet une unique solution, notée $\alpha$, dans l'intervalle [14~;~15].
		\item Recopier et compléter le tableau ci-dessous en faisant tourner étape par étape la fonction solution\_equation ci-contre, écrite en langage Python

\begin{minipage}{0.61\linewidth}
\begin{tabularx}{\linewidth}{|m{2.25cm}|c|*{4}{>{\centering \arraybackslash}X|}}\hline
	$a$		&14&&&&\\ \hline
	$b$		&15&&&&\\ \hline
	$b - a$	&1&&&&\\ \hline
	$m$		&14,5&&&&\cellcolor[gray]{0.2}\\ \hline
Condition $f(m) > 8$	&FAUX	&	&	&	&\cellcolor[gray]{0.2}\\ \hline
\end{tabularx}
\end{minipage}\hfill
\begin{minipage}{0.36\linewidth}
{\footnotesize\begin{tabular}{|l|}\hline
from math import exp \\
def f(x):\\
\qquad return (20*x+8)*exp(-1/4*x)\\
~\\
def solution\_equation():\\
\qquad a,b = 14,15\\
\qquad while b-a > 0.1 :\\
\qquad\begin{tabular}{|l}
\quad m = (a+b)/2\\
\quad if f(m) > 8:\\
\quad|\quad a = m\\
\quad else:\\
\quad |\quad  b = m\\
\end{tabular}\\
\qquad return a,b\\ \hline
\end{tabular}}

\end{minipage}
		\item Quel est l'objectif de la fonction solution\_equation
dans le contexte de la question ?
	\end{enumerate}
\end{enumerate}

\bigskip
